% ============================================================
% NDH Direction Geometry
% Module: 06-direction-geometry.tex
% Altitude: A4–A7 (Governance → Stability → Direction → Identity)
% Membrane: Direction-Constitution • Apex-Gradient Field
% ============================================================

\section{Direction Geometry}
\label{sec:direction-geometry}

The Direction Constitution governs the preferred tangent field of the NDH
manifold. It defines the apex-gradient field, regulates directional flow,
and constrains how geometric and operator-level structures may evolve along
constitutionally permitted directions. Direction operators preserve identity,
respect curvature bounds, and interact with the Quantum Sponge to ensure
non-absorptive directional reconstruction.

% ------------------------------------------------------------
% 1. Preferred Tangent Field
% ------------------------------------------------------------

\subsection{Preferred Tangent Field}

Let $T(X)$ denote the tangent bundle of the NDH manifold $X$. The Direction
Constitution defines a preferred tangent field

\[
\vec{d} : X \to T(X)
\]

assigning to each point $x \in X$ a direction vector $\vec{d}(x)$.

The preferred tangent field governs constitutional directionality.

% ------------------------------------------------------------
% 2. Apex-Gradient Field
% ------------------------------------------------------------

\subsection{Apex-Gradient Field}

The apex-gradient field is defined as the gradient of a constitutional
potential function $\Phi : X \to \mathbb{R}$:

\[
\nabla \Phi(x) = \vec{d}(x).
\]

The potential $\Phi$ is constitutionally constrained:

\[
\Phi(x) \le \Phi_{\max}.
\]

The apex-gradient field determines the constitutional flow of geometric
structures.

% ------------------------------------------------------------
% 3. Directional Flow
% ------------------------------------------------------------

\subsection{Directional Flow}

Directional flow is defined as the integral curve of the preferred tangent
field:

\[
\gamma'(t) = \vec{d}(\gamma(t)).
\]

Directional flow must satisfy:

\begin{itemize}
    \item \textbf{Curvature safety}:
    \[
    \kappa(\gamma(t)) \le \kappa_{\max}.
    \]

    \item \textbf{Identity preservation}:
    \[
    \pi(\gamma(t)) = \mathrm{id}_{\mathcal{M}}.
    \]

    \item \textbf{Non-absorptive evolution}:
    \[
    \gamma(t_1) = \gamma(t_2) \Rightarrow t_1 = t_2.
    \]
\end{itemize}

Directional flow integrates with holonomy stability.

% ------------------------------------------------------------
% 4. Direction Operators
% ------------------------------------------------------------

\subsection{Direction Operators}

Direction operators act on the preferred tangent field:

\[
\mathcal{O}_{\mathrm{dir}} :
T(X) \to T(X).
\]

They must satisfy:

\begin{itemize}
    \item \textbf{Apex-gradient preservation}:
    \[
    \mathcal{O}_{\mathrm{dir}}(\vec{d}) = \vec{d}.
    \]

    \item \textbf{Curvature envelope preservation}:
    \[
    \kappa'(x) \le \kappa_{\max}.
    \]

    \item \textbf{Identity preservation}:
    \[
    \pi \circ \mathcal{O}_{\mathrm{dir}} = \pi.
    \]
\end{itemize}

Direction operators govern constitutional directionality under $\mathcal{C}_5$.

% ------------------------------------------------------------
% 5. Interaction with Adjacency Categories
% ------------------------------------------------------------

\subsection{Adjacency Interaction}

Direction geometry interacts with adjacency categories via directional lifts:

\[
\mathcal{D}_s : \mathcal{A}_s \to T(X), \qquad
\mathcal{D}_c : \mathcal{A}_c \to T(X), \qquad
\mathcal{D}_o : \mathcal{A}_o \to T(X).
\]

These lifts ensure that spectral, constitutional, and operator-level
adjacencies respect constitutional directionality.

% ------------------------------------------------------------
% 6. Interaction with Quantum Sponge
% ------------------------------------------------------------

\subsection{Quantum Sponge Interaction}

The Quantum Sponge interacts with direction geometry via the functorial lift

\[
\mathcal{Q}_{\mathrm{dir}} :
T(X) \to \mathbf{Cat}.
\]

This ensures that reconstruction behavior remains directionally safe.

% ------------------------------------------------------------
% 7. Constitutional Integration
% ------------------------------------------------------------

\subsection{Constitutional Integration}

The Direction Constitution is represented by the functor $\mathcal{C}_5$:

\[
\mathcal{C}_5 : \mathcal{M} \to \mathcal{M}.
\]

It must satisfy:

\[
\mathcal{C}_5(\vec{d}) = \vec{d}, \qquad
\mathcal{C}_5(\Phi) \le \Phi_{\max}.
\]

This integrates direction geometry with the constitutional functor tower.

% ------------------------------------------------------------
% End of Module
% ------------------------------------------------------------
